\PassOptionsToPackage{unicode}{hyperref}
\PassOptionsToPackage{hyphens}{url}
\documentclass[11pt,a4paper]{article}

\usepackage[T1]{fontenc}
\usepackage[utf8]{inputenc}
\usepackage[brazil]{babel}
\usepackage{lmodern}
\usepackage{microtype}

\usepackage{xcolor}
\definecolor{javaOrange}{HTML}{E76F00}
\definecolor{javaDeep}{HTML}{BF5700}
\definecolor{javaBlue}{HTML}{1565C0}
\definecolor{javaTeal}{HTML}{00695C}
\definecolor{javaAmber}{HTML}{B45309}
\definecolor{javaInk}{HTML}{2B2140}
\definecolor{javaCodeBg}{HTML}{FFF8F0}
\definecolor{javaRule}{HTML}{FFD7B5}
\definecolor{javaShade}{HTML}{FFF3E8}

\usepackage{tcolorbox}
\tcbuselibrary{skins,breakable}

\newtcbox{\inlinecode}{
  on line,
  colback=javaCodeBg, colframe=javaRule,
  boxrule=0.5pt, arc=2pt,
  fontupper=\ttfamily\small\color{javaDeep},
  boxsep=1pt, left=2pt, right=2pt, top=1pt, bottom=1pt,
}

\usepackage[margin=2.4cm, top=2.8cm, bottom=2.8cm]{geometry}

\usepackage{fancyhdr}
\pagestyle{fancy}
\fancyhf{}
\renewcommand{\headrulewidth}{0.4pt}
\renewcommand{\headrule}{\hbox to\headwidth{\color{javaRule}\leaders\hrule height \headrulewidth\hfill}}
\fancyhead[L]{\small\sffamily\color{javaDeep} Java Programming MOOC}
\fancyhead[R]{\small\sffamily\color{javaDeep} Apresentação do Curso}
\fancyfoot[C]{\small\sffamily\color{gray}\thepage}

\usepackage{titlesec}
\titleformat{\section}
  {\sffamily\Large\bfseries\color{javaOrange}}{\thesection}{0.7em}{}
  [\vspace{2pt}{\color{javaRule}\titlerule[1.2pt]}]
\titleformat{\subsection}
  {\sffamily\large\bfseries\color{javaDeep}}{\thesubsection}{0.6em}{}
\titleformat{\subsubsection}
  {\sffamily\normalsize\bfseries\color{javaTeal}}{\thesubsubsection}{0.5em}{}

\usepackage{enumitem}
\setlist[itemize]{itemsep=2pt, topsep=4pt, parsep=0pt}
\setlist[enumerate]{itemsep=2pt, topsep=4pt, parsep=0pt}

\usepackage{booktabs}
\usepackage{array}
\usepackage{colortbl}
\usepackage{longtable}

\usepackage{hyperref}
\hypersetup{
  colorlinks=true,
  linkcolor=javaBlue, urlcolor=javaBlue, citecolor=javaBlue,
  pdftitle={Programação em Java — Apresentação do Curso},
  pdfauthor={Java Programming MOOC · Universidade de Helsinque},
  pdflang={pt-BR},
}

\newtcolorbox{destaque}[1][Atenção]{
  enhanced, breakable,
  colback=javaShade, colframe=javaOrange,
  leftrule=4pt, rightrule=0.4pt, toprule=0.4pt, bottomrule=0.4pt,
  arc=2pt,
  fonttitle=\sffamily\bfseries\color{javaOrange},
  title=#1,
  top=4pt, bottom=4pt,
}

\newtcolorbox{faqbox}{
  enhanced, breakable,
  colback=white, colframe=javaRule,
  leftrule=3pt, rightrule=0.4pt, toprule=0.4pt, bottomrule=0.4pt,
  arc=2pt,
  top=4pt, bottom=4pt, left=8pt, right=6pt,
}

\makeatletter
\newcommand{\subtitle}[1]{\gdef\@subtitle{#1}}
\renewcommand{\maketitle}{%
  \begin{center}
    {\color{javaRule}\rule{\linewidth}{1pt}}\\[6pt]
    {\Huge\sffamily\bfseries\color{javaOrange} \@title}\\[6pt]
    \ifx\@subtitle\undefined\else
      {\Large\sffamily\color{javaDeep} \@subtitle}\\[4pt]
    \fi
    {\small\sffamily\color{gray} \@author}\\[2pt]
    {\small\sffamily\color{gray} \@date}\\[6pt]
    {\color{javaRule}\rule{\linewidth}{1pt}}
  \end{center}
  \vspace{1em}
}
\makeatother

\title{Programação em Java}
\subtitle{Apresentação do Curso: Sobre, Suporte e Perguntas Frequentes}
\author{Java Programming MOOC · Universidade de Helsinque — tradução para o português}
\date{2026}

\begin{document}
\maketitle
\tableofcontents
\vspace{1.5em}

% ─────────────────────────────────────────────────────────────────────────────
\section{Sobre o Curso}

Este é um curso \textbf{gratuito e online} de programação em Java oferecido pela \textbf{Universidade de Helsinque}, desenvolvido pelo grupo de pesquisa \href{https://www.helsinki.fi/en/researchgroups/data-driven-education}{Agile Education Research}. Trata-se de uma versão atualizada do antigo curso \textit{Object-Oriented Programming with Java}.

O objetivo é ensinar os \textbf{fundamentos de programação de computadores}, incluindo algoritmos e programação orientada a objetos, utilizando a linguagem \textbf{Java}.

\subsection{Para quem é este curso?}

O curso \textbf{não exige conhecimento prévio em programação}. É voltado tanto para iniciantes absolutos quanto para quem deseja aprender Java especificamente.

\subsection{Estrutura}

\begin{center}
\begin{tabular}{l c p{8cm}}
\toprule
\textbf{Curso} & \textbf{Partes} & \textbf{Conteúdo} \\
\midrule
Java Programming I  & 1 a 7  & Fundamentos: variáveis, condicionais, repetição, métodos, listas, OOP, algoritmos \\
Java Programming II & 8 a 14 & Avançado: herança, polimorfismo, interfaces, generics, streams, lambdas \\
\bottomrule
\end{tabular}
\end{center}

\subsection{O que você aprenderá}

Ao concluir o \textbf{Java Programming I}, você será capaz de:

\begin{itemize}
  \item Escrever programas estruturados em Java
  \item Usar variáveis, condicionais e laços de repetição
  \item Criar e chamar métodos
  \item Trabalhar com listas, arrays e Strings
  \item Projetar e implementar classes e objetos
  \item Aplicar encapsulamento e composição
  \item Implementar algoritmos básicos de ordenação e busca
  \item Entender os paradigmas de programação
\end{itemize}

\subsection{Ferramentas necessárias}

\begin{enumerate}
  \item \textbf{Java Development Kit (JDK) versão 11 ou superior}\\
    \url{https://www.oracle.com/java/technologies/downloads/}\\
    Alternativa gratuita: \url{https://adoptium.net/}

  \item \textbf{IntelliJ IDEA} (IDE recomendada)\\
    \url{https://www.jetbrains.com/idea/download/}

  \item \textbf{Plugin TMC (Test My Code)}\\
    \url{https://www.mooc.fi/en/installation/intellij/}
\end{enumerate}

\subsection{Como começar}

\begin{enumerate}
  \item Acesse \url{https://www.mooc.fi/en/sign-up} e crie uma conta gratuita
  \item Instale o JDK 11 e o IntelliJ IDEA
  \item Instale o plugin TMC no IntelliJ
  \item No TMC, selecione o curso \textbf{"Java Programming I"}
  \item Faça o download da Parte 1 e comece!
\end{enumerate}

\subsection{Links principais}

\begin{center}
\begin{tabular}{l l}
\toprule
\textbf{Recurso} & \textbf{Link} \\
\midrule
Página principal      & \url{https://java-programming.mooc.fi/} \\
Criar conta           & \url{https://www.mooc.fi/en/sign-up} \\
Ver seus pontos       & \url{https://tmc.mooc.fi} \\
Ver certificados      & \url{https://www.mooc.fi/en/profile/completions} \\
Repositório GitHub    & \url{https://github.com/rage/java-programming} \\
Guia de instalação    & \url{https://www.mooc.fi/en/installation/intellij/} \\
\bottomrule
\end{tabular}
\end{center}

% ─────────────────────────────────────────────────────────────────────────────
\section{Avaliação e Certificados}

A avaliação é feita exclusivamente por meio dos \textbf{exercícios de programação}, corrigidos automaticamente pela ferramenta \textbf{Test My Code (TMC)}.

Para o \textbf{certificado de conclusão}, você precisa atingir pelo menos \textbf{80\% dos pontos} em cada parte:

\begin{itemize}
  \item Certificado do \textbf{Java Programming I}: ao completar 80\%+ das partes 1 a 7
  \item Certificado do \textbf{Java Programming II}: ao completar 80\%+ das partes 8 a 14
\end{itemize}

Certificados disponíveis em: \url{https://www.mooc.fi/en/profile/completions}

\begin{destaque}[Atenção]
Este curso \textbf{não oferece créditos acadêmicos} da Universidade de Helsinque. Para cursos com créditos ECTS: \url{https://programming-23.mooc.fi/}.
\end{destaque}

% ─────────────────────────────────────────────────────────────────────────────
\section{Suporte e Assistência}

\subsection{Canal oficial: Discord}

O suporte é oferecido \textbf{voluntariamente} no servidor Discord do Departamento de Ciência da Computação da Universidade de Helsinque.

\textbf{Entre no servidor:}
\url{https://study.cs.helsinki.fi/discord/join/java}

\subsection{Boas práticas ao pedir ajuda}

Ao postar uma dúvida, inclua:

\begin{enumerate}
  \item \textbf{Qual exercício} você está tentando resolver
  \item \textbf{O que você tentou fazer} — mostre seu código
  \item \textbf{Qual erro aparece} — copie a mensagem completa
  \item \textbf{O que você esperava} que acontecesse
\end{enumerate}

\subsection{Links de apoio ao estudo}

\begin{center}
\begin{tabular}{l l}
\toprule
\textbf{Recurso} & \textbf{Link} \\
\midrule
Discord do curso        & \url{https://study.cs.helsinki.fi/discord/join/java} \\
Repositório GitHub      & \url{https://github.com/rage/java-programming} \\
Documentação Java       & \url{https://docs.oracle.com/en/java/javase/11/docs/api/} \\
Tutorial Java Oracle    & \url{https://docs.oracle.com/javase/tutorial/} \\
W3Schools Java          & \url{https://www.w3schools.com/java/} \\
Stack Overflow          & \url{https://stackoverflow.com/questions/tagged/java} \\
Visualizador Java       & \url{https://pythontutor.com/java.html} \\
IntelliJ IDEA           & \url{https://www.jetbrains.com/idea/download/} \\
Guia TMC                & \url{https://www.mooc.fi/en/installation/intellij/} \\
\bottomrule
\end{tabular}
\end{center}

% ─────────────────────────────────────────────────────────────────────────────
\section{Perguntas Frequentes (FAQ)}

\subsection{Sobre o curso em geral}

\begin{faqbox}
\textbf{O curso é gratuito?}\\
Sim. O curso é completamente gratuito — material, exercícios e certificado.
\end{faqbox}

\vspace{0.5em}
\begin{faqbox}
\textbf{Preciso criar uma conta?}\\
Sim. Para usar o TMC, crie uma conta em \url{https://www.mooc.fi/en/sign-up}.
\end{faqbox}

\vspace{0.5em}
\begin{faqbox}
\textbf{Posso estudar no meu próprio ritmo?}\\
Sim. Não há prazos fixos para conclusão.
\end{faqbox}

\vspace{0.5em}
\begin{faqbox}
\textbf{Este curso oferece créditos acadêmicos?}\\
Não. Apenas certificado de conclusão. Para créditos ECTS: \url{https://programming-23.mooc.fi/}.
\end{faqbox}

\subsection{Sobre instalação e ferramentas}

\begin{faqbox}
\textbf{Qual IDE devo usar?}\\
IntelliJ IDEA (Community Edition, gratuita) com o plugin TMC.\\
Instruções: \url{https://www.mooc.fi/en/installation/intellij/}
\end{faqbox}

\vspace{0.5em}
\begin{faqbox}
\textbf{Qual versão do Java devo instalar?}\\
JDK versão 11. Remova outras versões para evitar conflitos.
\end{faqbox}

\vspace{0.5em}
\begin{faqbox}
\textbf{Estou recebendo erro sobre JAVA\_HOME. Como resolver?}\\
Verifique se o JDK 11 está instalado e configure a variável de ambiente \inlinecode{JAVA\_HOME} apontando para o diretório de instalação do JDK.
\end{faqbox}

\subsection{Sobre o plugin TMC e exercícios}

\begin{faqbox}
\textbf{Como atualizo o plugin TMC?}\\
No IntelliJ IDEA: \textbf{Help → Check for Updates}.
\end{faqbox}

\vspace{0.5em}
\begin{faqbox}
\textbf{Como baixo os exercícios das próximas partes?}\\
O TMC libera automaticamente ao atingir \textbf{25\% dos pontos} da parte anterior.
\end{faqbox}

\vspace{0.5em}
\begin{faqbox}
\textbf{Como reporto um possível bug em um exercício?}\\
Primeiro verifique se o erro está no seu código (peça ajuda no Discord). Se confirmar que é um bug, notifique os assistentes no Discord.
\end{faqbox}

% ─────────────────────────────────────────────────────────────────────────────
\section{Créditos}

O curso foi criado e é mantido pelo grupo \textbf{Agile Education Research} da Universidade de Helsinque:

\begin{itemize}
  \item \textbf{Arto Hellas} e \textbf{Matti Luukkainen} (autores principais)
  \item Colaboradores: Antti Laaksonen, Henrik Nygren, Juho Leinonen, entre outros
\end{itemize}

\textbf{Licença:} Creative Commons BY-NC-SA 4.0\\
\url{https://creativecommons.org/licenses/by-nc-sa/4.0/}

\begin{center}
\begin{tabular}{l l}
\toprule
\textbf{Recurso} & \textbf{Link} \\
\midrule
Universidade de Helsinque   & \url{https://www.helsinki.fi/en} \\
Agile Education Research    & \url{https://www.helsinki.fi/en/researchgroups/data-driven-education} \\
Plataforma mooc.fi          & \url{https://www.mooc.fi/en/} \\
Test My Code (TMC)          & \url{https://tmc.mooc.fi} \\
GitHub do curso             & \url{https://github.com/rage/java-programming} \\
\bottomrule
\end{tabular}
\end{center}

\end{document}
